\chapter{The Kitaev Chain as the Fermionic Image of the XY Mother Model}
\label{chap:kitaev-chain-fermionic-image}

The Kitaev chain is the fermionic image of the transverse-field XY chain.  From
the spin point of view it is obtained by the Jordan--Wigner transformation; from
the fermion point of view it is a one-dimensional spinless $p$-wave
superconductor.  This chapter isolates the fermionic interpretation of the XY
solution and connects it to Majorana edge modes and the BdG topological
invariant.

\section{Spinless \texorpdfstring{$p$}{p}-wave superconducting chain}
\label{sec:kitaev-p-wave-chain}

We use the convention
\begin{equation}
\label{eq:kitaev-hamiltonian}
    H_{\rm K}
    =
    -t\sum_{j}
    \left(
        c_j^\dagger c_{j+1}+c_{j+1}^\dagger c_j
    \right)
    -\mu\sum_j\left(n_j-\frac12\right)
    +\Delta\sum_j
    \left(
        c_jc_{j+1}+c_{j+1}^\dagger c_j^\dagger
    \right),
\end{equation}
where $n_j=c_j^\dagger c_j$.  The hopping $t$ and chemical potential $\mu$ are
ordinary number-conserving terms, while $\Delta$ is a nearest-neighbor odd-parity
pairing amplitude.  For the moment $\Delta$ is taken real.  A uniform complex
phase will be discussed below.

The pairing term breaks fermion-number conservation but preserves fermion parity:
\begin{equation}
\label{eq:kitaev-parity-conservation}
    [H_{\rm K},(-1)^N]=0,
    \qquad
    N=\sum_j n_j.
\end{equation}
Thus the Kitaev chain belongs naturally to the same BdG framework used for the
XY chain.

\section{Dictionary from the XY mother model}
\label{sec:kitaev-xy-dictionary}

The bulk Jordan--Wigner image of the XY chain in the conventions of these notes
is
\begin{equation}
\label{eq:kitaev-xy-bulk-image}
    H_{\rm XY}^{\rm bulk}
    =
    -J\sum_j
    \left(
        c_j^\dagger c_{j+1}+c_{j+1}^\dagger c_j
    \right)
    -J\gamma\sum_j
    \left(
        c_j^\dagger c_{j+1}^\dagger+c_{j+1}c_j
    \right)
    -h\sum_j(1-2n_j).
\end{equation}
Using
\begin{equation}
    c_{j+1}^\dagger c_j^\dagger=-c_j^\dagger c_{j+1}^\dagger,
    \qquad
    c_jc_{j+1}=-c_{j+1}c_j,
\end{equation}
comparison between \cref{eq:kitaev-hamiltonian} and
\cref{eq:kitaev-xy-bulk-image} gives the bulk dictionary
\begin{equation}
\label{eq:kitaev-xy-dictionary}
    t=J,
    \qquad
    \Delta=J\gamma,
    \qquad
    \mu=-2h.
\end{equation}
The constant term matches because
\begin{equation}
    -\mu\sum_j\left(n_j-\frac12\right)
    =
    2h\sum_j n_j-hL
    \qquad
    (\mu=-2h).
\end{equation}

This dictionary is convention-dependent.  If one writes the pairing term as
$\Delta_{\rm cr}\sum_j(c_j^\dagger c_{j+1}^\dagger+c_{j+1}c_j)$ instead, then
$\Delta_{\rm cr}=-J\gamma$ in the present Jordan--Wigner convention.  The
spectrum and phase diagram are unchanged; only the orientation convention for
the BdG pairing plane is changed.

For a periodic spin chain, the boundary term is also sector dependent.  In a
fixed fermion-parity sector,
\begin{equation}
\label{eq:kitaev-boundary-from-spin}
    P_F=p
    \quad\Longrightarrow\quad
    c_{L+1}=-p\,c_1.
\end{equation}
For an intrinsically fermionic Kitaev chain, one may instead specify periodic,
anti-periodic, or open boundary conditions directly.  The bulk topological phase
diagram is the same, but exact finite-size spectra depend on this choice.

\section{BdG block and spectrum}
\label{sec:kitaev-bdg-block-spectrum}

For a translation-invariant chain, the normal-state dispersion is
\begin{equation}
\label{eq:kitaev-xi-k}
    \xi(k)=-\mu-2t\cos k.
\end{equation}
With the Fourier and BdG convention used in the earlier chapters, the odd-parity
pairing gap for real $\Delta$ may be written as
\begin{equation}
\label{eq:kitaev-gap-real}
    \Delta(k)=2\ii\Delta\sin k.
\end{equation}
The BdG block is therefore
\begin{equation}
\label{eq:kitaev-bdg-block}
    \mathcal{H}_{\rm K}(k)
    =
    \begin{pmatrix}
        -\mu-2t\cos k & 2\ii\Delta\sin k\\
        -2\ii\Delta\sin k & \mu+2t\cos k
    \end{pmatrix}.
\end{equation}
Equivalently,
\begin{equation}
\label{eq:kitaev-d-vector}
    \mathcal{H}_{\rm K}(k)
    =
    \bm d_{\rm K}(k)\cdot\bm\tau,
    \qquad
    \bm d_{\rm K}(k)
    =
    \bigl(0,-2\Delta\sin k,-\mu-2t\cos k\bigr).
\end{equation}
The quasiparticle spectrum is
\begin{equation}
\label{eq:kitaev-spectrum}
    E(k)
    =
    \sqrt{(-\mu-2t\cos k)^2+4\Delta^2\sin^2 k}.
\end{equation}
Under \cref{eq:kitaev-xy-dictionary}, this reproduces the XY-chain spectrum
\begin{equation}
    E(k)
    =
    2\sqrt{(h-J\cos k)^2+(J\gamma\sin k)^2}.
\end{equation}

For $\Delta\neq0$, the bulk gap can close only at $k=0$ or $k=\pi$.  The two
conditions are
\begin{equation}
\label{eq:kitaev-gap-closing}
    k=0:\quad \mu=-2t,
    \qquad
    k=\pi:
    \quad \mu=2t.
\end{equation}
Thus the gapped phases are separated by the lines
\begin{equation}
\label{eq:kitaev-critical-lines}
    \mu=\pm 2t.
\end{equation}

\section{Complex pairing phase and Nambu gauge}
\label{sec:kitaev-complex-gap-phase}

Let the pairing amplitude be complex:
\begin{equation}
\label{eq:kitaev-complex-delta}
    \Delta=\abs{\Delta}\ee^{\ii\theta_\Delta}.
\end{equation}
With the Fourier convention containing a constant phase $\phi$,
\begin{equation}
    c_k=\frac{\ee^{-\ii\phi}}{\sqrt{L}}
    \sum_j\ee^{-\ii kj}c_j,
\end{equation}
the momentum-space gap becomes
\begin{equation}
\label{eq:kitaev-complex-gap-fourier-phase}
    \Delta_\phi(k)
    =
    2\ii\abs{\Delta}\,
    \ee^{\ii(\theta_\Delta-2\phi)}\sin k.
\end{equation}
The phase $\phi$ is a Nambu/Fourier gauge convention.  Choosing
\begin{equation}
\label{eq:kitaev-gauge-remove-uniform-phase}
    \phi=\frac{\theta_\Delta}{2}
\end{equation}
removes the uniform superconducting phase from the bulk BdG block.  Therefore a
spatially uniform pairing phase is not by itself a new bulk phase of matter.  It
rotates the BdG vector in the pairing plane but leaves $E(k)$ and the topological
criterion unchanged.

A phase difference, phase winding, or boundary flux is different.  If the pairing
phase cannot be removed by a single-valued global gauge transformation, it may
represent physical supercurrent or flux data and can affect boundary spectra.

\section{Majorana representation}
\label{sec:kitaev-majorana-representation}

Define site Majorana operators using the convention of the covariance-matrix
chapter:
\begin{equation}
\label{eq:kitaev-majorana-definitions}
    A_j=c_j+c_j^\dagger,
    \qquad
    B_j=\ii(c_j-c_j^\dagger).
\end{equation}
They satisfy
\begin{equation}
    A_j^\dagger=A_j,
    \qquad
    B_j^\dagger=B_j,
    \qquad
    \{A_i,A_j\}=2\delta_{ij},
    \qquad
    \{B_i,B_j\}=2\delta_{ij},
    \qquad
    \{A_i,B_j\}=0.
\end{equation}
Also
\begin{equation}
\label{eq:kitaev-density-majorana}
    n_j-\frac12=-\frac{\ii}{2}A_jB_j.
\end{equation}
For real $\Delta$, the Hamiltonian becomes
\begin{equation}
\label{eq:kitaev-majorana-hamiltonian}
    H_{\rm K}
    =
    \frac{\ii}{2}\sum_j
    \left[
        \mu A_jB_j
        +(t-\Delta)A_jB_{j+1}
        -(t+\Delta)B_jA_{j+1}
    \right],
\end{equation}
up to boundary terms fixed by the chosen boundary condition.

Two limiting cases reveal the phase structure immediately.

\paragraph{Trivial atomic limit.}
For $t=\Delta=0$ and $\mu<0$ or $\mu>0$, Majoranas on the same site are paired by
$\ii\mu A_jB_j/2$.  There are no protected edge Majoranas.

\paragraph{Topological dimerized limit.}
For $\mu=0$ and $\Delta=t>0$,
\begin{equation}
\label{eq:kitaev-topological-special-point}
    H_{\rm K}
    =
    -\ii t\sum_j B_jA_{j+1}.
\end{equation}
On an open chain this pairs $B_j$ with $A_{j+1}$ on each bond.  The operators
$A_1$ and $B_L$ are left unpaired and become zero-energy Majorana edge modes.
The two edge Majoranas combine into a nonlocal fermion
\begin{equation}
\label{eq:kitaev-edge-fermion}
    f_{\rm edge}=\frac{1}{2}(A_1+\ii B_L),
\end{equation}
producing two nearly degenerate ground states of opposite fermion parity in a
long open chain.

\section{Topological invariant}
\label{sec:kitaev-topological-invariant}

For real parameters, the BdG vector lies in a plane.  Define
\begin{equation}
\label{eq:kitaev-q-function}
    q(k)
    =
    d_z(k)+\ii d_y(k)
    =
    -\mu-2t\cos k-2\ii\Delta\sin k.
\end{equation}
The integer winding number in the real-pairing BDI convention is
\begin{equation}
\label{eq:kitaev-winding-number}
    \nu
    =
    \frac{1}{2\pi\ii}
    \int_{-\pi}^{\pi}\dd k\,\partial_k\log q(k).
\end{equation}
For $t\Delta\neq0$, the loop winds around the origin if and only if
\begin{equation}
\label{eq:kitaev-topological-condition}
    \abs{\mu}<2\abs{t}.
\end{equation}
The orientation of the winding depends on the sign of $t\Delta$ and on the
chosen definition of $q(k)$, but the distinction between zero and nonzero winding
is robust.

If only particle--hole symmetry is kept, the class-D invariant is the
corresponding $\ZZ_2$ quantity.  In the present convention it can be written as
\begin{equation}
\label{eq:kitaev-pfaffian-invariant}
    \mathcal{M}
    =
    \sgn\{\xi(0)\xi(\pi)\}
    =
    \sgn\{(-\mu-2t)(-\mu+2t)\}
    =
    \sgn(\mu^2-4t^2).
\end{equation}
The topological phase has
\begin{equation}
\label{eq:kitaev-class-d-nontrivial}
    \mathcal{M}=-1
    \quad\Longleftrightarrow\quad
    \abs{\mu}<2\abs{t}.
\end{equation}
Using the XY dictionary $\mu=-2h$ and $t=J$, this becomes
\begin{equation}
\label{eq:kitaev-xy-topological-condition}
    \abs{h}<\abs{J},
\end{equation}
which is the region between the two Ising gap-closing lines of the XY mother
model.

\section{Relation to the TFIM and the XY chain}
\label{sec:kitaev-relation-to-tfim-xy}

The transverse-field Ising model corresponds to the Kitaev chain with
\begin{equation}
\label{eq:kitaev-tfim-dictionary}
    \gamma=1
    \quad\Longleftrightarrow\quad
    \Delta=t=J,
    \qquad
    \mu=-2h.
\end{equation}
For $J,h>0$, the TFIM transition at $h=J$ is the Kitaev transition at
\begin{equation}
    \mu=-2t.
\end{equation}
The spin ferromagnetic phase $h<J$ is the topological superconducting phase
$\abs{\mu}<2\abs{t}$.  The spin paramagnet $h>J$ is the trivial superconducting
phase.

This identification is conceptually useful but should not be overinterpreted.
The fermions in the Kitaev-chain description are Jordan--Wigner fermions, not
microscopic electrons of the original spin chain.  The Majorana edge modes of
the open Kitaev chain correspond, in the spin language, to the near degeneracy of
the symmetry-broken Ising phase and to nonlocal string operators.

\section{Summary}
\label{sec:kitaev-summary}

The Kitaev chain is the fermionic form of the XY mother model.  The convention
used in these notes gives the bulk dictionary
\begin{equation}
    t=J,
    \qquad
    \Delta=J\gamma,
    \qquad
    \mu=-2h.
\end{equation}
Its BdG spectrum is
\begin{equation}
    E(k)=\sqrt{(-\mu-2t\cos k)^2+4\abs{\Delta}^2\sin^2 k}.
\end{equation}
For nonzero pairing, the bulk transitions occur at $\mu=\pm2t$.  The phase
$\abs{\mu}<2\abs{t}$ is topological and supports Majorana edge modes on an open
chain.  A uniform complex phase of $\Delta$ rotates the BdG pairing plane and can
be absorbed by a Nambu gauge choice, while phase winding or boundary flux may
carry physical information.
